\subsection{Credit Assignment in RLVR}
Reinforcement learning with verifiable rewards (RLVR) methods such as GRPO~\citep{grpo} have become a standard paradigm for post-training reasoning models because they optimize against automatically checkable outcomes without human preference labels~\citep{grpo,deepseekr1}. Their central limitation is that the verifier provides a sequence-level signal, so every token in a rollout inherits the same advantage regardless of whether it corresponds to a pivotal deduction or a stylistic filler. This makes long-horizon reasoning fundamentally a credit assignment problem. A common response is to introduce process reward models (PRMs) or step-level value estimators that score intermediate reasoning steps~\citep{lightman2024lets,wang2024math,yang2025testtimepromptintervention,luo2024improve,chen2024step,zhang2024generative,yang2025dynamicearlyexitreasoning,dai2025sgrpoearlyexitreinforcement,cui2025process}. Some of these methods rely on costly human step annotations, while others replace them with automated or implicit supervision; in either case, they require auxiliary reward modeling and extra computation beyond the base policy.

More recent work seeks fine-grained credit assignment while staying strictly within a verifier-only RLVR pipeline. These methods reshape token-level updates using model-internal proxies such as entropy, uncertainty, key-token statistics, attention dynamics, or outcome sensitivity~\citep{cheng2026reasoning,wangbeyond,chen2025seed,sun2025ktaemodelfreealgorithmkeytokens,xie2025unlocking,li2025attention,chen2025beyond,li2026outcome}. While they show that uniform sequence-level advantages are unnecessarily coarse, they estimate token importance entirely from intrinsic heuristics. Our method is complementary in motivation yet fundamentally different in mechanism: rather than relying on heuristic proxies, we extract rigorous token-level assessments via self-distillation under privileged context, while keeping the update direction anchored to the verifier reward.

\subsection{On-Policy Distillation}
Distillation has also been explored as an alternative source of dense token-level supervision for reasoning models. In on-policy distillation (OPD), a stronger teacher evaluates the student's own rollouts and provides token-level targets along on-policy trajectories~\citep{agarwal2024gkd,lu2025onpolicy,fu2026revisitingonpolicydistillationempirical}; recent reports suggest that strong teacher-guided distillation can even rival or complement RL-based post-training~\citep{mimo2025flash}. However, standard OPD requires maintaining an external, typically larger teacher throughout training. To eliminate this overhead, recent self-distillation methods let the same model act as both student and teacher, with the teacher conditioned on privileged information such as reference solutions or verifier feedback~\citep{zhao2026opsd,hbotter2026reinforcementlearningselfdistillation_sdpo}. Closely related on-policy distillation variants have also been explored for continual learning from demonstrations, context internalization, and reasoning compression, where the teacher is conditioned on demonstrations, auxiliary context, or conciseness instructions~\citep{shenfeld2026selfdistillation,ye2026policy,sang2026policy,penaloza2026privileged}. A recent concurrent work, TRRD \citep{zhang2026reinforcementawareknowledgedistillationllm}, also identifies the conflict between additive KL penalties and reward maximization, proposing to inject teacher probabilities into the policy importance ratio. However, TRRD operates on the trust region anchor while RLSD directly modulates the advantage magnitude via a Bayesian evidence ratio.

The shared design choice across these methods is to retain a distribution-matching objective between a privileged or context-conditioned teacher and a less-informed student. Our work starts from a similar setting but departs from that paradigm. We show that when the teacher's prediction depends on latent information unavailable to the student, matching the teacher distribution becomes structurally ill-posed and can induce privileged information leakage and unstable long-run training. RLSD therefore does not use the teacher as a generative target; instead, it repurposes the privileged discrepancy only as a scalar multiplier for credit assignment, decoupling dense token-level magnitude from the optimization direction. 
